% Intended LaTeX compiler: pdflatex
\documentclass[a4paper,12pt]{article}
\usepackage[utf8]{inputenc}
\usepackage[T1]{fontenc}
\usepackage{graphicx}
\usepackage{longtable}
\usepackage{wrapfig}
\usepackage{rotating}
\usepackage[normalem]{ulem}
\usepackage{amsmath}
\usepackage{amssymb}
\usepackage{capt-of}
\usepackage{hyperref}
\usepackage{\string~/.emacs.d/common}
\author{mahmood}
\date{\today}
\title{rule of inference}
\hypersetup{
 pdfauthor={mahmood},
 pdftitle={rule of inference},
 pdfkeywords={},
 pdfsubject={},
 pdfcreator={Emacs 28.2 (Org mode 9.5.5)}, 
 pdflang={English}}
\begin{document}

\maketitle
\begin{note}
this document and others can be found on my website\\
\url{https://mahmoodsheikh36.github.io/}
\end{note}
see \url{https://en.wikipedia.org/wiki/List\_of\_rules\_of\_inference}\\
\begin{definition}
a \textbf{rule of inference} is a logical form consisting of a function which takes \href{20230524212241-premise.org}{premise}s, analyzes their syntax, and returns a \href{20230524212250-conclusion.org}{conclusion} (or conclusions).\\
\end{definition}
rules of inference are usually given in the following standard form:\\
\[
\infer{conclusion}{premise_1 & premise_2 & \cdots & premise_n}
\]\\
\begin{characteristic}
a rule of inference is sound if the \textbf{conclusion} logically follows from the \textbf{premises}\\
\begin{my_example}
the following isnt a sound rule of inference\\
\[
\infer{\beta}{\alpha \lor \beta}
\]\\
the following is a sound rule of inference\\
\[
\infer{\beta}{\alpha \land \beta}
\]\\
\end{my_example}
\end{characteristic}
we arrive from premises to a conclusion with the \href{20230525071323-formal_proof.org}{inference process}
\end{document}